\section[Material Platform: Ge-on-Si]{Material and Integration Platform: Germanium on Silicon}
\label{sec:pd_material}

With the waveguide-integrated PIN structure identified as the correct device class
and the O-band wavelength fixed at \SI{1310}{\nano\metre}, the absorber material
must now be chosen. The selection follows directly from the absorption physics of
\cref{sec:pd_fundamentals}: only germanium combines strong O-band absorption,
sufficient carrier velocity, and silicon-foundry process compatibility. This
section elaborates the three dimensions of that argument, namely spectral
suitability, CMOS integration, and material-induced trade-offs, at the depth
required to motivate the design choices of \cref{sec:pd_design}.

\subsection[Why Germanium at 1310 nm]{Why Germanium at 1310 nm}
\label{subsec:pd_bandstructure}

Germanium is an indirect-gap semiconductor whose lowest conduction-band minimum
lies at the $L$-point of the Brillouin zone, separated from the valence-band
maximum at $\Gamma$ by an indirect gap
$E_{g,\mathrm{ind}}=\SI{0.66}{\electronvolt}$ at \SI{300}{\kelvin}. A secondary
conduction-band minimum at the $\Gamma$-point defines a direct gap
$E_{g,\mathrm{dir}}=\SI{0.80}{\electronvolt}$~\cite{AshcroftMermin2021,Rossler2009,Ziman1972}.
The photon energy at \SI{1310}{\nano\metre} is
$\hbar\omega=hc/\lambda\approx\SI{0.946}{\electronvolt}$, which exceeds
$E_{g,\mathrm{dir}}$ by approximately \SI{0.15}{\electronvolt}. Optical absorption
at this wavelength therefore proceeds predominantly through \emph{direct} interband
transitions at the $\Gamma$-point, for which crystal momentum is intrinsically
conserved without phonon assistance~\cite{Basu2003,Kantorovich2004}. Silicon,
whose lowest direct transition lies at \SI{3.4}{\electronvolt} and whose indirect
gap of \SI{1.12}{\electronvolt} makes it effectively transparent at O-band
wavelengths, cannot be used as the absorber. This single material-physics argument
is the entire justification for the Ge-on-Si heterostructure.

\paragraph{Biaxial strain and the direct-gap red-shift.}
The application of biaxial tensile strain $\varepsilon$ reduces the direct gap
according to the linear deformation-potential relation:
\begin{equation}
    E_{g,\mathrm{dir}}(\varepsilon)
    \approx E_{g,\mathrm{dir}}(0)
      - \bigl(2a_v + a_c\bigr)\bigl(1 - c_{12}/c_{11}\bigr)\varepsilon,
    \label{eq:bandgap_strain}
\end{equation}
where $a_v$ and $a_c$ are the valence- and conduction-band hydrostatic
deformation potentials and $c_{11}$, $c_{12}$ are the elastic stiffness
constants. The biaxial tensile strain of approximately
\SIrange{0.2}{0.3}{\percent} that develops in Ge epitaxial layers during
post-growth cooldown, owing to the thermal expansion coefficient mismatch
between Ge and Si, reduces $E_{g,\mathrm{dir}}$ by approximately
\SI{25}{\milli\electronvolt} relative to the unstrained bulk value, red-shifting
the onset of direct-gap absorption by roughly \SI{50}{\nano\metre} and enhancing
$\kappa(\lambda)$ at C-band wavelengths. For the O-band operating point of the
present device this correction is small but it is the principal mechanism
exploited in strained-Ge extended-wavelength detector designs.

\begin{table}[H]
\centering
\caption{Material properties of candidate photodetector absorbers at
\SI{300}{\kelvin}. $\alpha$ values are referenced to \SI{1310}{\nano\metre}.
The final column summarises compatibility with a standard silicon-photonics
process flow.}
\label{tab:material_compare}
\footnotesize
\begin{tabular*}{\linewidth}{@{\extracolsep{\fill}}L{2.8cm}C{1.6cm}C{2.0cm}C{2.2cm}C{2.0cm}C{1.2cm}L{2.3cm}@{}}
\toprule
Material & $E_g$ [eV] & $\alpha_{1310}$ [\si{\per\centi\metre}] &
$v_{\mathrm{sat},h}$ [\si{\centi\metre\per\second}] &
$\mu_p$ [\si{\centi\metre\squared\per\volt\per\second}] &
$\Delta a/a_{\mathrm{Si}}$ [\%] & CMOS compatibility \\
\midrule
Ge & 0.66 (ind.), 0.80 (dir.) & $\sim2100$ & \SI{4e6}{} & 1900 & 4.2 & Compatible \\
Si & 1.12 (ind.) & $\sim10$ & \SI{2.5e6}{} & 480 & 0 & Compatible \\
InGaAs (lattice-matched to InP) & 0.75 & $\sim1\times10^4$ & \SI{5e6}{} & 4500 & $>5$ & Not compatible \\
Ge$_{0.92}$Sn$_{0.08}$ & $\sim0.55$ (dir.) & $>5000$ & $\sim$\SI{4e6}{} & $\sim$1500 & 4.8 & Research-stage \\
\bottomrule
\end{tabular*}
\end{table}

\begin{figure}[H]
    \centering
    \missingfigure[figwidth=0.88\linewidth]{
        Two-panel concept figure. Left panel: absorption coefficient
        $\alpha(\lambda)$ versus wavelength for Ge, Si, InGaAs, and
        GeSn, with markers at \SI{1310}{\nano\metre} and
        \SI{1550}{\nano\metre}. Right panel: schematic comparison of the Ge
        and Si band structures showing the O-band photon energy relative to
        the Ge direct gap and the much larger Si direct-transition energy.}
    \caption{Material selection argument for O-band detection on a silicon
    photonics platform. Germanium combines appreciable absorption at
    \SI{1310}{\nano\metre}, near-direct-gap behaviour, high hole saturation
    velocity, and foundry-compatible heterogeneous integration.}
    \label{fig:ge_vs_si_material}
\end{figure}

\subsection[CMOS Compatibility and Integration Constraints]{CMOS Compatibility and Integration Constraints}
\label{subsec:pd_ge_integration}

Germanium can be grown selectively on silicon by chemical vapour deposition (CVD)
using a two-step heteroepitaxial process: a low-temperature ($\sim\SI{400}{\celsius}$)
nucleation layer of \SIrange{30}{50}{\nano\metre} seeds the growth, and a
high-temperature ($\sim\SI{650}{\celsius}$) bulk growth achieves the target
thickness. Cyclic thermal annealing between \SIrange{600}{750}{\celsius}
reduces the threading dislocation density from
$\sim\SI{1e9}{\per\centi\metre\squared}$ after growth to
$\sim\SI{5e7}{\per\centi\metre\squared}$ after annealing, which is the practical
lower bound achievable without compositionally graded buffer layers. The total
thermal budget of this process is compatible with the front-end-of-line (FEOL)
constraints of a \SI{220}{\nano\metre} SOI platform.

Major silicon-photonics foundries, including Intel, imec, and GlobalFoundries,
have qualified Ge-on-Si photodetector processes at \SI{220}{\nano\metre} SOI node.
The process compatibility removes the III-V bonding step required by InGaAs
receivers, eliminates the need for flip-chip assembly, and enables monolithic
co-integration of the photodetector with Si waveguides, grating couplers,
multiplexers, and electro-optic modulators on the same die.

\paragraph{Heterojunction band alignment and its numerical consequences.}
The electrical consequence of the Ge/Si material interface directly conditions
the boundary conditions of the drift-diffusion problem. Owing to the lattice
mismatch ($\Delta a/a\approx4.2\%$) and the difference in band alignment, the
Ge/Si heterojunction sustains a valence-band discontinuity of
$\Delta E_v\approx\SI{0.54}{\electronvolt}$ and a conduction-band offset
$\Delta E_c\approx\SI{-0.14}{\electronvolt}$~\cite{Basu2003}. For a $p^{++}$-Si
anode contacting the intrinsic Ge absorber, the large $\Delta E_v$ acts as a
potential barrier suppressing hole injection from silicon into germanium, thereby
reducing the minority-carrier diffusion contribution to dark current quantified by
\cref{eq:jdark_diff}. A thin compositionally graded transition layer of thickness
$\delta_{\mathrm{grad}}\approx\SI{5}{\nano\metre}$ at the heterojunction removes
the valence-band singularity that would otherwise prevent convergence of the
self-consistent drift-diffusion solver without perceptibly modifying the
electrostatic potential profile in the intrinsic germanium bulk.

\subsection[Material-Induced Trade-offs]{Material-Induced Trade-offs}
\label{subsec:pd_ge_tradeoffs}

The same lattice mismatch ($\Delta a/a\approx4.2\%$) that makes Ge-on-Si
heteroepitaxy necessary also creates a residual defect population that cannot be
completely eliminated by annealing. Threading dislocations at density
TDD $\approx\SI{5e7}{\per\centi\metre\squared}$ act as Shockley-Read-Hall
recombination centres in the Ge bulk, directly setting the SRH lifetime $\tau_0$
and the dark current floor through \cref{eq:jdark_srh}. The current CHARGE
implementation uses $\tau_0=\SI{1}{\nano\second}$ for Ge, which represents a
conservative nanosecond-scale recombination environment for defect-limited
Ge-on-Si material.

The interaction between Ge thickness, dislocation density, and device performance
creates a three-way trade-off. Thicker Ge layers absorb more light, improving
$\mathcal{A}$ via \cref{eq:abs_fraction}, but slow carrier transit through
\cref{eq:bw_tt} and increase the total SRH generation volume, raising $I_d$ via
\cref{eq:jdark_srh}. Thinner Ge layers improve speed and reduce dark current but
sacrifice absorption efficiency. Additionally, the near-interface Ge layer,
specifically the first \SIrange{50}{100}{\nano\metre} grown on Si, has higher
defect density than the bulk, which is why the $n^{++}$-Ge cathode of the
present device is limited to \SI{50}{\nano\metre}: it minimises both free-carrier absorption (FCA)
in the contact region and the volume of highest-defect-density material within
the optical generation zone.

A compact material comparison is therefore sufficient to close the absorber
selection argument: Ge is not chosen because it is universally superior to all
alternatives, but because it is the only material in \cref{tab:material_compare}
that simultaneously satisfies the optical, transport, and integration constraints
of the present O-band receiver.

\subsection[Carrier Transport and Dark Current Mechanisms]{Carrier Transport Physics and Dark Current Mechanisms in Ge-on-Si Photodiodes}
\label{subsec:pd_transport_dark}

The performance of a waveguide-integrated Ge-on-Si photodiode is governed not
only by the optical absorption physics described in
\cref{subsec:pd_bandstructure}, but equally by the dynamics of photogenerated
carrier transport and the unavoidable generation of thermally excited carriers
that constitute the device dark current. These two phenomena are deeply coupled:
the same electric field that sweeps photogenerated carriers across the intrinsic
germanium absorber with high velocity also governs the rate at which
trap-mediated thermal generation replenishes the depletion region with
spurious carriers. Understanding both mechanisms from first principles is
essential before the geometric and doping trade-offs of the device design in
\cref{sec:pd_design} can be motivated rigorously.

\subsubsection[Drift, Diffusion, and Carrier Saturation Velocity]{Drift, Diffusion, and the Onset of Velocity Saturation}
\label{subsubsec:transport_drift}

Under the applied reverse bias $V_R$ that fully depletes the intrinsic
germanium layer of thickness $w_i$, the internal electric field $\mathcal{E}$
in the depletion region drives photogenerated electrons and holes toward the
cathode and anode, respectively, by drift. At moderate fields, the mean drift
velocity of a carrier species $s \in \{e, h\}$ is proportional to
$\mathcal{E}$ through the low-field mobility $\mu_s$:
\begin{equation}
    v_{d,s} = \mu_s \mathcal{E}.
    \label{eq:drift_lowfield}
\end{equation}
For bulk germanium at \SI{300}{\kelvin}, $\mu_e \approx \SI{3900}{\centi\metre\squared\per\volt\per\second}$
and $\mu_h \approx \SI{1900}{\centi\metre\squared\per\volt\per\second}$,
both substantially exceeding their silicon counterparts. The high hole mobility
of Ge is of particular practical relevance in a $p$-$i$-$n$ structure where
hole transit across the full intrinsic thickness can otherwise constitute the
bandwidth-limiting transit-time contribution.

As $\mathcal{E}$ increases beyond approximately
\SI{10}{\kilo\volt\per\centi\metre} for electrons and
\SI{5}{\kilo\volt\per\centi\metre} for holes in germanium, inter-valley and
optical-phonon scattering rates grow faster than the energy gained from the
field, causing the velocity to saturate asymptotically. The empirical
Canali model provides a compact and physically consistent description of
velocity saturation in semiconductor materials:
\begin{equation}
    v_{d,s}(\mathcal{E})
    = \frac{\mu_s \mathcal{E}}
      {\Bigl[1 + \bigl(\mu_s \mathcal{E} / v_{\mathrm{sat},s}\bigr)^{\beta_s}\Bigr]^{1/\beta_s}},
    \label{eq:canali_velocity}
\end{equation}
where $v_{\mathrm{sat},s}$ is the saturation velocity of species $s$ and
$\beta_s$ is a dimensionless exponent that interpolates the transition between
the linear and saturated regimes. For holes in germanium,
$v_{\mathrm{sat},h} \approx \SI{4e6}{\centi\metre\per\second}$ and
$\beta_h \approx 1$, which reduces \cref{eq:canali_velocity} to the
hyperbolic form familiar from standard textbook treatments. In the
high-field limit relevant to a reverse-biased photodiode, \cref{eq:canali_velocity}
collapses to $v_{d,s} \to v_{\mathrm{sat},s}$, independently of mobility.
This saturation regime is precisely the operating point targeted by the
present device, since a sufficiently large reverse bias ensures that carrier
velocities are clamped at $v_{\mathrm{sat}}$ and are therefore independent
of any spatial non-uniformity in $\mathcal{E}$ along the propagation axis of
the optical mode.

\paragraph{Transit-time limited bandwidth.}
The direct device-level consequence of carrier velocity saturation is the
transit-time contribution to the $3$\,dB bandwidth. If holes traverse the
intrinsic germanium thickness $w_i$ at their saturated velocity, the
transit-time-limited bandwidth is:
\begin{equation}
    f_{\mathrm{tt}}
    = \frac{0.45\, v_{\mathrm{sat},h}}{w_i},
    \label{eq:mat_bw_tt}
\end{equation}
where the numerical prefactor 0.45 follows from the Fourier analysis of the
induced current waveform in the Shockley--Ramo framework for a uniformly
distributed generation profile. \Cref{eq:bw_tt} establishes a direct
inverse relationship between absorber thickness and transit-time bandwidth:
halving $w_i$ doubles $f_{\mathrm{tt}}$, but at the cost of a proportional
reduction in the absorbed fraction of the evanescently coupled optical power,
as governed by the Beer-Lambert relation discussed in \cref{sec:pd_fundamentals}.
This thickness trade-off is the central geometric tension of the design
and motivates the use of waveguide-coupling geometry in preference to
surface illumination: by distributing absorption along the propagation
direction rather than the vertical direction, the coupling length $L_{\mathrm{pd}}$
independently controls the absorbed fraction while $w_i$ remains free to be
minimised for bandwidth.

\subsubsection[Dark Current: Physical Origins and Analytical Decomposition]{Dark Current: Physical Origins and Analytical Decomposition}
\label{subsubsec:dark_current}

In the absence of illumination, a reverse-biased germanium photodiode passes a
current $I_d$ that constitutes the fundamental noise floor of the receiver.
Dark current enters the signal chain as a shot-noise source with
single-sided power spectral density $S_I = 2q I_d$, which adds in quadrature
with the signal shot noise and the amplifier noise floor. Minimising $I_d$
is therefore not merely a power-efficiency consideration but a direct
bandwidth-noise trade-off: at the data rates of 400GBASE-DR4, even
moderate dark currents of order \SI{1}{\micro\ampere} can degrade the
receiver sensitivity by several tenths of a decibel if the transimpedance
amplifier noise is not sufficiently dominant.

Three physically distinct mechanisms contribute to the total dark current
density $J_d$ in a reverse-biased Ge-on-Si $p$-$i$-$n$ diode:
diffusion of minority carriers from the quasi-neutral contact regions,
Shockley--Read--Hall (SRH) thermal generation within the depletion region,
and, at sufficiently large reverse bias, band-to-band tunnelling enhanced
by the narrow direct gap of germanium. Each mechanism exhibits a characteristic
bias and temperature dependence that allows them to be distinguished
experimentally and that must be accounted for separately in the device model.

\paragraph{Diffusion current.}
Minority carriers generated thermally within one diffusion length of the
edge of the depletion region drift across the junction under the action of
the built-in and applied fields. For the $p^{++}$-Si anode contacting the
$i$-Ge absorber of the present device, the analysis is modified by the
Ge/Si heterojunction band alignment: the valence-band discontinuity
$\Delta E_v \approx \SI{0.54}{\electronvolt}$ at the Si/Ge interface acts as
a potential barrier that suppresses minority hole injection from the Si anode
into the Ge absorber, substantially reducing the silicon contribution to
diffusion dark current compared with a homojunction geometry. The residual
diffusion current density arising from minority electrons in the $p^{++}$-Ge
anode and minority holes in the $n^{++}$-Ge cathode takes the form:
\begin{equation}
    J_{\mathrm{diff}}
    = q\, n_i^2 \left(
        \frac{D_e}{L_e N_A} + \frac{D_h}{L_h N_D}
      \right),
    \label{eq:jdark_diff}
\end{equation}
where $n_i$ is the intrinsic carrier concentration of germanium,
$D_{e,h}$ are the minority carrier diffusion coefficients,
$L_{e,h} = \sqrt{D_{e,h}\tau_{e,h}}$ are the corresponding diffusion
lengths, and $N_A$, $N_D$ are the acceptor and donor concentrations in the
$p^{++}$ and $n^{++}$ contact regions, respectively. The strong exponential
temperature dependence of $n_i^2 \propto \exp(-E_g / k_B T)$ combined with
the relatively narrow gap of germanium ($E_g = \SI{0.66}{\electronvolt}$)
makes $J_{\mathrm{diff}}$ significantly larger in Ge devices than in
comparable Si junctions at the same temperature, and explains why Ge
photodiodes typically exhibit dark currents in the range
\SIrange{10}{100}{\nano\ampere} even at modest device areas.

\paragraph{Shockley--Read--Hall generation current.}
The dominant dark-current mechanism in the depletion region of a defect-rich
Ge-on-Si device is trap-assisted thermal generation via mid-gap states. In the
SRH framework~\cite{ShockleyRead1952,Hall1952}, a defect level at energy
$E_t$ near the intrinsic Fermi level mediates a two-step process by which
valence electrons are excited to the conduction band via the trap, generating
an electron-hole pair without the need for a direct band-to-band phonon or
photon interaction. The volumetric generation rate under full depletion
approximation is maximised when $E_t = E_i$ and reduces to:
\begin{equation}
    G_{\mathrm{SRH}}
    = \frac{n_i}{2\tau_0},
    \label{eq:srh_gen_rate}
\end{equation}
where $\tau_0 = (\tau_e \tau_h)^{1/2}$ is the geometric-mean SRH lifetime,
which encapsulates the product of the thermal capture cross-section $\sigma$,
the trap density $N_t$, and the thermal velocity $v_{\mathrm{th}}$ through
$\tau_0^{-1} = \sigma v_{\mathrm{th}} N_t$. Integrating $G_{\mathrm{SRH}}$
over the depleted volume $V_{\mathrm{dep}} = A \cdot w_i$, where $A$ is
the device cross-sectional area, yields the SRH dark current:
\begin{equation}
    I_{\mathrm{SRH}}
    = q \, G_{\mathrm{SRH}} \, V_{\mathrm{dep}}
    = \frac{q\, n_i\, A\, w_i}{2\tau_0}.
    \label{eq:jdark_srh}
\end{equation}

The physical significance of \cref{eq:jdark_srh} in the Ge-on-Si context
is immediate. Threading dislocations at density TDD $\approx \SI{5e7}{\per\centi\metre\squared}$,
which persist after cyclic thermal annealing as discussed in
\cref{subsec:pd_ge_tradeoffs}, introduce mid-gap states whose spatial
density $N_t$ is directly proportional to the TDD. Effective SRH lifetimes in
processed Ge-on-Si devices are commonly treated as nanosecond-scale rather than
bulk-germanium values ($\tau_0 \sim \SI{1}{\micro\second}$). The repository
therefore adopts $\tau_0=\SI{1}{\nano\second}$ as the Ge electrical boundary
condition in the drift-diffusion simulations of \cref{subsec:pd_charge}.

Three design implications follow directly from \cref{eq:jdark_srh}. First,
$I_{\mathrm{SRH}}$ scales linearly with the depleted volume $w_i \cdot A$,
providing a strong incentive to minimise both the absorber thickness and the
device footprint. Second, the linear dependence on $n_i$ means that
$I_{\mathrm{SRH}}$ is exponentially sensitive to temperature through
$n_i \propto \exp(-E_g / 2k_B T)$; this temperature scaling is less severe
than that of the diffusion component because it involves $n_i$ rather than
$n_i^2$, making SRH generation the dominant mechanism at lower temperatures
where the diffusion contribution collapses more rapidly. Third, because
$I_{\mathrm{SRH}}$ is proportional to $w_i$ whereas the absorbed optical
power scales as $1 - \exp(-\alpha w_i)$ for a surface-illuminated geometry,
there exists an optimal absorber thickness that maximises the ratio of
photocurrent to dark current; in the waveguide-coupled geometry this optimum
shifts, since the absorption is accumulated along the propagation axis and
$w_i$ can be reduced without sacrificing responsivity, as noted above.

\paragraph{Band-to-band tunnelling.}
At reverse biases beyond approximately \SIrange{2}{3}{\volt} for the doping
concentrations considered here, the electric field in the depletion region
becomes sufficiently large that quantum-mechanical tunnelling of valence
electrons directly into conduction-band states across the narrow direct gap
of germanium produces a field-dependent excess dark current. The
Zener tunnelling current density scales as~\cite{Sze2021}:
\begin{equation}
    J_{\mathrm{BTBT}}
    \propto \mathcal{E}^{1/2} \exp\!\left(
        -\frac{4\sqrt{2m^*}\, E_{g,\mathrm{dir}}^{3/2}}{3 q \hbar \mathcal{E}}
    \right),
    \label{eq:btbt}
\end{equation}
where $m^*$ is the reduced tunnelling effective mass and $\hbar$ is the
reduced Planck constant. The exponential dependence of $J_{\mathrm{BTBT}}$
on $\mathcal{E}^{-1}$ and on $E_{g,\mathrm{dir}}^{3/2}$ means that, for
the direct gap of $\SI{0.80}{\electronvolt}$ of germanium, tunnelling becomes
non-negligible at fields accessible under practical reverse-bias conditions.
Because the present device is intended to operate at reverse biases of
\SIrange{1}{2}{\volt}, the field in the intrinsic region remains well below
the Zener onset threshold and $J_{\mathrm{BTBT}}$ is negligible in the nominal
operating regime. Nevertheless, the proximity of the germanium direct gap to
the tunnelling onset motivates a careful upper bound on the reverse bias
used during high-temperature burn-in or reliability qualification, where
the thermally softened band structure and field crowding near device corners
could locally activate the tunnelling mechanism.

\subsubsection[Total Dark Current and the Shot Noise Floor]{Total Dark Current Budget and Its Consequence for Receiver Sensitivity}
\label{subsubsec:dark_noise}

Within the nominal operating regime, the total dark current is:
\begin{equation}
    I_d
    = I_{\mathrm{diff}} + I_{\mathrm{SRH}},
    \label{eq:total_idark}
\end{equation}
with $I_{\mathrm{diff}}$ governed by \cref{eq:jdark_diff} and $I_{\mathrm{SRH}}$
by \cref{eq:jdark_srh}. For a device of cross-section
$A \approx \SI{0.5}{\micro\metre} \times \SI{0.4}{\micro\metre}$
and intrinsic thickness $w_i \approx \SI{400}{\nano\metre}$,
evaluated at \SI{300}{\kelvin} with a nanosecond-scale $\tau_0$, the
SRH component is of order tens of nanoamperes, consistent with the
experimentally reported dark currents of \SIrange{5}{50}{\nano\ampere} for
similarly sized Ge-on-Si waveguide photodetectors in the literature.

The receiver-level impact is quantified through the dark-current shot noise
contribution to the total input-referred noise current spectral density. In
a transimpedance-amplifier front-end with bandwidth $B$, the root-mean-square
dark current noise is:
\begin{equation}
    i_{n,d} = \sqrt{2 q I_d B},
    \label{eq:shot_noise_dark}
\end{equation}
which for $I_d = \SI{20}{\nano\ampere}$ and $B = \SI{26.6}{\giga\hertz}$
(corresponding to a single $100$\,Gb/s lane of 400GBASE-DR4 at a $4:1$
PAM-4 symbol rate of $26.6$\,GBd) yields
$i_{n,d} \approx \SI{14}{\nano\ampere\per\sqrt{\hertz}} \times \sqrt{B}
\approx \SI{2.3}{\micro\ampere}_{\mathrm{rms}}$.
This figure is typically sub-dominant relative to the transimpedance amplifier
input-referred noise current in practical receivers, confirming that the dark
current of a well-optimised Ge-on-Si photodiode does not constitute the
sensitivity-limiting noise source at the data rates of interest. However, any
degradation of $\tau_0$ due to process variation or elevated operating
temperature will increase $I_{\mathrm{SRH}}$ and can shift this balance,
providing the physical rationale for the tight thermal-budget control and
cyclic annealing protocol discussed in \cref{subsec:pd_ge_integration}.

The analytical decomposition of \cref{eq:total_idark} thus closes the argument
connecting material defect physics, device geometry, and receiver noise
performance: the SRH lifetime $\tau_0$ set by the threading dislocation density
of the Ge-on-Si heteroepitaxial process propagates directly into a dark current
floor, which in turn determines the minimum shot-noise pedestal against which
optical signals must be detected. The design choices of absorber thickness,
contact doping concentration, and reverse bias voltage explored in
\cref{sec:pd_design} are all conditioned on minimising \cref{eq:total_idark}
subject to the competing constraints on absorption efficiency and bandwidth
established by \cref{eq:bw_tt} and the Beer-Lambert absorption analysis.
